% Scope: Connect relaxation models, exchange and susceptibility to sequence-dependent tissue contrast without treating tissue constants as universal.
% Status: architecture only; no chapter prose or completed problems yet.
% Prerequisites: Chapters 2, 3, 7 and 8.
% Planned worked examples: Competing tissue contrast and multi-pool relaxation.
% Planned figures: Spectral densities, exchange and contrast curves.
% Planned challenge problems: Test a monoexponential assumption; compare nulling under imperfect inversion.
\chapter{Contrast and Relaxation}
\label{ch:09}

\section{Microscopic origins of relaxation}
\label{sec:09:microscopic-origins-of-relaxation}
\subsection{Fluctuating fields and correlation functions}
\subsection{Spectral density, T1 and T2}
\subsection{Field strength and molecular-motion regimes}

\section{Susceptibility and transverse dephasing}
\label{sec:09:susceptibility-and-transverse-dephasing}
\subsection{Microscopic and macroscopic field variation}
\subsection{Reversible and irreversible signal loss}
\subsection{Spin-echo versus gradient-echo contrast}

\section{Exchange and magnetization transfer}
\label{sec:09:exchange-and-magnetization-transfer}
\subsection{Two-pool exchange models}
\subsection{Bound-pool saturation and MT contrast}
\subsection{MTR, quantitative MT and model dependence}

\section{Contrast preparation}
\label{sec:09:contrast-preparation}
\subsection{Inversion, saturation and T2 preparation}
\subsection{Fat suppression and spectral selectivity}
\subsection{Spin locking and T1-rho concepts}

\section{Contrast agents and tissue models}
\label{sec:09:contrast-agents-and-tissue-models}
\subsection{Relaxivity, concentration and nonlinear effects}
\subsection{Compartmentation and exchange limitations}
\subsection{Tissue variability and sequence dependence}

\section{Contrast interpretation}
\label{sec:09:contrast-interpretation}
\subsection{Weighting versus parameter measurement}
\subsection{Clinical uses of T1, T2 and susceptibility}
\subsection{Failure of simple contrast heuristics}

\section{Worked examples}
\label{sec:09:worked-examples}
% Planned calculations: Competing tissue contrast and multi-pool relaxation.
\subsection{Derivation and quantitative calculation}
\subsection{Assumptions, limiting cases and scanner interpretation}

\section{Challenge problems}
\label{sec:09:challenge-problems}
% Problem design: Test a monoexponential assumption; compare nulling under imperfect inversion.
% Use challengeproblem and stable labels prob:09:<topic>.
% Complete solutions belong in Appendix E, section for this chapter.
% Use \begin{solution}{prob:09:<topic>} ... \end{solution} there.
\subsection{Derivation and quantitative design}
\subsection{Model criticism and integrated reasoning}
